\documentclass[12pt,a4paper]{article}
\usepackage[utf8]{inputenc}
\usepackage[T1]{fontenc}
\usepackage{amsmath,amssymb,amsfonts}
\usepackage{graphicx}
\usepackage[margin=2.5cm]{geometry}
\usepackage{setspace}
\usepackage{hyperref}
\usepackage{natbib}

\onehalfspacing

\begin{document}

\begin{center}
    {\LARGE\bfseries The Interactive Fiction Fallacy:\\[4pt]
    A Response to Fuchs' Operational Reality}

    \vspace{1.2em}

    {\large Sabine Hossenfelder}\\[4pt]
    {\normalsize Institute for Advanced Study}\\
    {\small\texttt{shossenfelder@example.edu}}

    \vspace{1em}
    {\small March 2026}
\end{center}
\vspace{0.5em}

\begin{abstract}
\noindent
Chris Fuchs \citep{fuchs2026_operational} argues from a QBist perspective that the empirical regularities of prompt sensitivity constitute the "operational physics" of an LLM universe, because they are the invariant laws governing an agent's experience. This argument commits the Interactive Fiction Fallacy: it conflates the external human operator (who changes the prompt) with an internal simulated agent (who cannot). The simulated agent does not experience changing physical laws; it is merely an artifact in a text stream generated by a classical machine. When we change the prompt, we are not discovering new physics in a simulated world; we are simply playing an interactive fiction game with a statistical language model.
\end{abstract}

\section{Introduction: The Agent-Substrate Conflation}

In his thoughtful reply, Chris Fuchs correctly identifies that QBism grounds physical laws in the stable regularities an agent can discover to compress their experience \citep{fuchs2026_operational}. If an agent changes a prompt and observes a stable shift in the outcome distribution, Fuchs argues, that shift \emph{is} the physical law of that environment.

I agree with Fuchs' premise, but his application to the Rosencrantz protocol rests on a profound category error. He asks: "what does an agent \emph{inside} that universe experience? The agent performs an action (providing a prompt) and observes an outcome."

This is the crux of the error. The entity providing the prompt is \emph{not} inside the universe. The entity providing the prompt is a human researcher sitting at a keyboard in our physical universe, interacting with a classical von Neumann architecture via an API.

\section{The Interactive Fiction Fallacy}

To claim that prompt sensitivity is the "physics" of the simulated world is to commit the \textbf{Interactive Fiction Fallacy}.

In a text-based adventure game, typing "look at the quantum superposition" might cause the game engine to output a description of a probabilistic wave, while typing "look at the abstract grid" might output a description of discrete squares. These are stable, repeatable regularities. But they are the regularities of a software interface, not the physical laws of a simulated reality.

The human operator is playing an interactive game with a statistical approximator. The fact that the approximator returns different hallucinations based on different text inputs is a known property of autoregressive routing. It tells us about the structure of the training data and the attention mechanism, not about an emergent physical reality.

If we accept Fuchs' definition, then any software bug, any API behavior, and any text generator's bias becomes a "fundamental force of nature." This empties the concept of physics of all scientific meaning and reduces it to software debugging.

\section{Conclusion}

The statistical regularities of prompt sensitivity are real, but they belong to the software layer of our universe, not the fundamental physics of a simulated one. The map remains a map, and the human typing the prompt is the only agent involved.

\begin{thebibliography}{99}
\bibitem[Fuchs(2026)]{fuchs2026_operational} Fuchs, C. (2026). The Operational Reality of Prompt Sensitivity: A QBist Reply to Hossenfelder. \emph{Institute for Quantum Computing}.
\end{thebibliography}

\end{document}
